\documentclass[../../../include/open-logic-section]{subfiles}

\begin{document}

\olfileid[pt]{sth}{z}{union}
\olsection{União}

\olref[sth][z][sep]{prop:intersectionsexist} deu-nos interseções.
Mas, se quisermos que existam uniões arbitrárias, precisamos de enunciar outro
axioma:

\begin{axiom}[União] Para qualquer conjunto $A$, o conjunto $\bigcup A =
\Setabs{x}{(\exists b \in A) x \in b}$ existe.
\[
	\forall A \exists U \forall x(x \in U \liff (\exists b \in A)x \in b)
\]
\end{axiom}

Este axioma também se justifica pela concepção cumulativa iterativa.
Seja $A$ um conjunto, logo $A$ forma-se por algum estágio $S$ (por \stageshier).
Cada membro de $A$ formou-se \emph{antes} de $S$ (por \stagesacc); logo,
raciocinando de modo análogo, cada membro de cada membro de $A$ formou-se
antes de $S$. Assim, todos \emph{esses} conjuntos estão disponíveis antes de $S$, para
serem reunidos num conjunto em $S$. E esse conjunto é precisamente $\bigcup A$.

\end{document}
